\subsubsection{Formål med simulationsmotoren}

Forklar kort, at du har lavet en generel 2D PDE-engine, som kan simulere
forskellige PDE-systemer ved at indlæse et selvvalgt \texttt{rhs} og en
initialbetingelse gennem \texttt{PDESystem}.

\subsubsection{Diskretisering af domænet}

Beskriv gitteret:
\begin{itemize}
    \item Domæne: $[x_{\min}, x_{\max}] \times [y_{\min}, y_{\max}]$
    \item Opløsning: $n_x \times n_y$
    \item Gitterafstande: $\Delta x$ og $\Delta y$
\end{itemize}

\texttt{Grid2D} laver et uniformt gitter og bruger ghost-celler omkring det
fysiske domæne.

\subsubsection{Repræsentation af tilstanden}

Forklar, at tilstanden $u$ gemmes som enten ét felt eller flere koblede felter:
\[
    u(t,x,y) \in \mathbb{R}^{n_{\text{fields}}}.
\]
Motoren understøtter derfor både simple PDE'er og koblede PDE-systemer.

\subsubsection{Randbetingelser}

Beskriv, hvordan ghost-celler bruges til periodiske randbetingelser. I din
engine sættes værdierne på kanten ved at kopiere fra den modsatte side af
domænet.

\subsubsection{Numeriske differentialoperatorer}

Lav et afsnit om finite differences:
\begin{itemize}
    \item Første afledte: central differens
    \item Anden afledte: central differens
    \item Laplace-operator
    \item Advektion
\end{itemize}

Disse er implementeret i \texttt{num\_diff.py}.

\subsubsection{Tidsintegration}

Beskriv RK4-metoden som tidsintegrator:
\[
    u^{n+1} = u^n + \frac{\Delta t}{6}(k_1 + 2k_2 + 2k_3 + k_4).
\]
Din engine bruger præallokerede arrays til $k_1$, $k_2$, $k_3$ og $k_4$, så
der undgås unødvendige allokeringer under simulationen.

\subsubsection{Adaptiv tidsstyring}

Forklar, at tidssteppet justeres ud fra størrelsen af højresiden:
\[
    \frac{\partial u}{\partial t} = F(t,u).
\]
Hvis $|F|$ bliver stor, reduceres $\Delta t$; hvis systemet ændrer sig
langsomt, kan $\Delta t$ vokse.

\subsubsection{Simulationsloop}

Beskriv den samlede proces:
\begin{enumerate}
    \item PDE-systemet indlæses dynamisk.
    \item Gitteret oprettes.
    \item Initialbetingelsen beregnes.
    \item RK4 med adaptiv $\Delta t$ bruges til at udvikle systemet.
    \item Resultater gemmes med et fast \texttt{save\_dt}.
\end{enumerate}

Dette svarer til strukturen i \texttt{simulate\_system}.

\subsubsection{Datalagring}

Forklar, at snapshots gemmes løbende i en memmap og derefter pakkes i en
\texttt{.npz}-fil sammen med tidspunkter og konfigurationsparametre.

\subsubsection{Stabilitet og fejlhåndtering}

Beskriv kort:
\begin{itemize}
    \item Kontrol for \texttt{inf}
    \item Clipping af løsningen
    \item Begrænsning af $\Delta t$ mellem \texttt{min\_dt} og
          \texttt{max\_dt}
\end{itemize}
